\documentclass[12pt]{article}
\usepackage{amsmath,amssymb}
\usepackage{amsthm}
\usepackage{hyperref}
\usepackage{geometry}
\usepackage{booktabs}
\geometry{margin=1in}

\title{A Curvature Identity for the Fibonacci Digital-Root Cycle Modulo~9}

\author{Guy Rinzema\\
Independent researcher\\
\texttt{guy.rinzema@icloud.com}}

\date{June 2026}

%--- theorem environments
\newtheorem{thm}{Theorem}
\newtheorem{prop}{Proposition}
\theoremstyle{remark}
\newtheorem{rem}{Remark}

\begin{document}
\maketitle

\begin{abstract}
For each seed we form the $24$-periodic orbit of the Fibonacci recurrence modulo~$9$,
write its terms with the digital-root convention $0\mapsto 9$, and measure its
\emph{digital curvature} $S=\sum_n \kappa_n^2$, the sum of squared discrete second
differences. We make two points. First, on the $72$ full-period seeds modulo~$9$ the
curvature takes exactly three values, $\{648,972,1296\}$, each on $24$ seeds; the maximum
$1296$ is attained precisely by the classical Fibonacci cycle $(1,1)$ and its rotations.
This is a finite statement and we prove it by exhaustive enumeration. Second, and more
usefully, $S$ is the squared norm of a circulant operator and therefore has an exact
Fourier form, $S=\frac{16}{N}\sum_{k}\sin^4(\pi k/N)\,|\hat a_k|^2$. This identity
explains the classification: the curvature is a high-pass energy, and the Fibonacci orbit
maximizes it because it is the most high-frequency of the three orbit classes modulo~$9$,
carrying about $98\%$ of its curvature energy in the near-Nyquist half of the spectrum. We use the same lens to
describe how far the maximality property extends to other moduli (it is a small-modulus
phenomenon, not a law) and to retire a conjecture from earlier versions of this note.
This version corrects errors in Versions~1--3; see the revision history in
Section~\ref{sec:errata}.
\end{abstract}

\noindent\textbf{Keywords:} Fibonacci numbers; Pisano period; digital root;
discrete curvature; circulant operator; discrete Fourier transform.

\section{Introduction}

Let $F(n)$ be the Fibonacci numbers, $F(0)=0$, $F(1)=1$. The sequence $F(n)\bmod 9$ has
Pisano period $24$~\cite{Wall1960}. Replacing each residue $0$ by $9$ (the digital-root
convention used in base~$10$, since the digital root of an integer equals its value
modulo~$9$ with $0\mapsto 9$) yields the familiar cycle
\[
1,1,2,3,5,8,4,3,7,1,8,9,8,8,7,6,4,1,5,6,2,8,1,9.
\]
This $24$-tuple is well known, but its higher-order difference statistics are not usually
examined. We attach to every modulo-$9$ orbit a single number---its \emph{digital
curvature}, the energy of its discrete second difference---and ask what values that number
takes.

The answer modulo~$9$ is a clean trichotomy (Theorem~\ref{thm:trichotomy}). The point of
this version of the note, however, is not the trichotomy itself but the reason behind it.
The curvature functional is the squared norm of the cyclic second-difference operator,
which is circulant and hence diagonalized by the discrete Fourier transform. This gives an
exact spectral formula for $S$ (Proposition~\ref{prop:spectral}) from which the
classification, the identity of the maximizer, and the limits of that maximality all
follow transparently. The spectral identity is a standard fact of harmonic analysis on the
cycle; what is (to our knowledge) not previously recorded is its application to the
digital-root orbits of a Pisano period, and the resulting fingerprint of the classical
Fibonacci cycle.

We have tried to be careful about what is and is not established. The trichotomy is a
theorem proved by finite computation. The spectral identity is a theorem. The statement
that the Fibonacci orbit maximizes curvature is true for small moduli but provably fails
for larger ones, and we report it as the bounded empirical fact that it is. Earlier
versions of this note advanced a closed-form conjecture and a symmetry argument that turn
out to be, respectively, a numerical coincidence and an error; both are corrected below.

\section{Definitions}\label{sec:defs}

Fix a modulus $m$ (we take $m=9$ unless stated otherwise) and let $P=\pi(m)$ be its
Pisano period. For a seed $(a_0,a_1)\in(\mathbb Z/m\mathbb Z)^2$ generate the orbit
\[
a_{n+1}=a_n+a_{n-1}\pmod m .
\]
We restrict attention to seeds whose orbit has \emph{minimal period exactly $P$}
(``full-period'' seeds); these are the orbits that realize the full Pisano cycle. The
\emph{rotation class} of such a seed is the set of $P$ pairs $(a_n,a_{n+1})$ occurring
along its orbit---equivalently, its $P$ cyclic shifts. Display
each residue with
\[
d(r)=\begin{cases}m & r\equiv 0,\\ r & r\in\{1,\dots,m-1\},\end{cases}
\]
so displayed values lie in $\{1,\dots,m\}$, and set
\[
\kappa_n=d(a_{n+1})-2\,d(a_n)+d(a_{n-1})\in\mathbb Z,
\qquad
S=\sum_{n=0}^{P-1}\kappa_n^{2},
\]
all indices taken cyclically modulo $P$. We call $S$ the \emph{digital curvature} of the
orbit. Two remarks are worth making at the outset.

\begin{rem}[Convention dependence]
$S$ depends on the display convention $0\mapsto m$. With residues taken in
$\{0,\dots,m-1\}$ instead, the same orbit yields different second differences and a
different $S$. The quantity studied here is therefore an invariant of the
\emph{digital-root presentation} of the orbit, not of the abstract residue sequence. This
is deliberate---the digital-root cycle is the object of interest---but it should be stated
plainly.
\end{rem}

\begin{rem}[Curvature is a cyclic energy]
$\kappa_n=(D d)_n$ where $D$ is the cyclic second-difference operator, so
$S=\lVert Dd\rVert^2$ is the discrete Dirichlet-type energy of the displayed orbit on the
cycle $C_P$. This observation drives everything in Section~\ref{sec:spectral}.
\end{rem}

\section{The Trichotomy Modulo 9}\label{sec:trichotomy}

\begin{thm}\label{thm:trichotomy}
Among the $72$ full-period seeds modulo~$9$, the digital curvature takes exactly the three
values
\[
S\in\{648,\;972,\;1296\},
\]
each attained by exactly $24$ seeds (equivalently, by one full-period orbit apiece).
Moreover $S=1296$ if and only if the seed lies in the rotation class of $(1,1)$, and then
\[
S=1296=\sum_{n=0}^{23}\kappa_n^2 .
\]
\end{thm}

\begin{proof}
There are $81$ ordered seeds $(a_0,a_1)\in\{1,\dots,9\}^2$; of these, $72$ generate orbits
of minimal period $24$ and $9$ do not. Evaluating $S$ on each full-period orbit is a
finite computation, carried out in the script of Section~\ref{sec:verify}; it returns the
multiset $\{648{:}24,\,972{:}24,\,1296{:}24\}$, with the value $1296$ occurring exactly on
the $24$ rotations of the $(1,1)$ orbit. A finite exhaustive check is a complete proof of a
statement quantified over a finite set.
\end{proof}

We stress that the enumeration \emph{is} the proof. Earlier versions of this note tried to
derive the three values from a congruence argument; that argument is vacuous (all three
values are divisible by $9$, so reduction modulo~$9$ cannot distinguish them) and has been
removed. What survives from it is a correct structural observation about \emph{which}
orbits carry the three values.

\begin{rem}[The three values are the three unit classes]\label{rem:units}
The unit group $(\mathbb Z/9\mathbb Z)^\times=\{1,2,4,5,7,8\}$ acts on orbits by
$a_n\mapsto g\,a_n$. Because $F_{n+12}\equiv -F_n\pmod 9$ \emph{at the level of residues},
negation $g=-1\equiv 8$ maps the $(1,1)$ orbit to its own shift by $12$; hence the subgroup
$\{\pm1\}$ fixes each orbit and the six units collapse to the three classes
\[
\{\pm1\}=\{1,8\},\quad \{\pm2\}=\{2,7\},\quad \{\pm4\}=\{4,5\},
\]
i.e.\ $(\mathbb Z/9\mathbb Z)^\times/\{\pm1\}$. Scaling the $(1,1)$ orbit by a
representative of each class gives precisely the three curvature values:
\[
g\in\{1,8\}\!\to\!1296,\qquad g\in\{2,7\}\!\to\!972,\qquad g\in\{4,5\}\!\to\!648,
\]
verified directly. So the trichotomy is the statement that curvature is an invariant of
$(\mathbb Z/9\mathbb Z)^\times/\{\pm1\}$, taking distinct values on its three classes.
(This residue-level symmetry is what makes the unit-class action well defined; it is only
the antisymmetry of the displayed \emph{second differences} that the $0\mapsto 9$
convention breaks, as Section~\ref{sec:halves} shows.)
\end{rem}

Table~\ref{tab:kappa} records the second differences of the maximal $(1,1)$ orbit, for use
below. One checks directly that $\sum_n\kappa_n=0$ and $\sum_n\kappa_n^2=1296$.

\begin{table}[h!]
\centering
\[
\begin{array}{c|cccccccccccc}
n & 0 & 1 & 2 & 3 & 4 & 5 & 6 & 7 & 8 & 9 & 10 & 11\\\hline
a_n & 1 & 1 & 2 & 3 & 5 & 8 & 4 & 3 & 7 & 1 & 8 & 9\\
\kappa_n & 8 & 1 & 0 & 1 & 1 & -7 & 3 & 5 & -10 & 13 & -6 & -2
\end{array}
\]
\[
\begin{array}{c|cccccccccccc}
n & 12 & 13 & 14 & 15 & 16 & 17 & 18 & 19 & 20 & 21 & 22 & 23\\\hline
a_n & 8 & 8 & 7 & 6 & 4 & 1 & 5 & 6 & 2 & 8 & 1 & 9\\
\kappa_n & 1 & -1 & 0 & -1 & -1 & 7 & -3 & -5 & 10 & -13 & 15 & -16
\end{array}
\]
\caption{Second differences for the $(1,1)$ orbit. $\sum\kappa_n=0$, $\sum\kappa_n^2=1296$.}
\label{tab:kappa}
\end{table}

\section{Curvature as Spectral Energy}\label{sec:spectral}

The cyclic second-difference operator $D$ on length-$N$ sequences is circulant, hence
diagonalized by the discrete Fourier transform. This yields an exact formula for $S$.

\begin{prop}\label{prop:spectral}
Let $d=(d_0,\dots,d_{N-1})$ be the displayed orbit (of length $N=P$), with unnormalized DFT
$\hat d_k=\sum_{n}d_n e^{-2\pi i kn/N}$. Then
\[
S=\frac1N\sum_{k=0}^{N-1}\mu_k^2\,|\hat d_k|^2,
\qquad
\mu_k=2\cos\!\frac{2\pi k}{N}-2=-4\sin^2\!\frac{\pi k}{N},
\]
equivalently
\[
\boxed{\;S=\frac{16}{N}\sum_{k=0}^{N-1}\sin^4\!\Big(\frac{\pi k}{N}\Big)\,|\hat d_k|^2.\;}
\]
\end{prop}

\begin{proof}
$\kappa_n=(Dd)_n=d_{n+1}-2d_n+d_{n-1}$ with indices modulo $N$, so $D$ is the circulant
matrix with symbol $e^{2\pi i k/N}-2+e^{-2\pi i k/N}=2\cos(2\pi k/N)-2=\mu_k$. The Fourier
modes $e^{2\pi i kn/N}$ are its eigenvectors with eigenvalues $\mu_k$, so
$\widehat{Dd}_k=\mu_k\hat d_k$. By Parseval, $S=\lVert Dd\rVert^2=\frac1N\sum_k|\mu_k\hat
d_k|^2=\frac1N\sum_k\mu_k^2|\hat d_k|^2$, and $\mu_k^2=16\sin^4(\pi k/N)$.
\end{proof}

This is standard harmonic analysis on the cycle; we record it only because it converts the
opaque integer $S$ into something interpretable. The weight $16\sin^4(\pi k/N)$ vanishes at
the constant mode $k=0$, grows monotonically across $k=1,\dots,\lfloor N/2\rfloor$, and is
maximal at the Nyquist mode $k=N/2$. Thus \emph{digital curvature is a high-pass energy}:
it ignores the orbit's mean and slow variation and rewards rapid, near-alternating
oscillation. Call the modes with $\sin^2(\pi k/N)\ge\frac12$---equivalently
$N/4\le k\le 3N/4$---the \emph{near-Nyquist half} of the spectrum.

For the maximal $(1,1)$ orbit ($N=24$, $S=1296$), Table~\ref{tab:spectrum} gives the
per-mode contribution $c_k=\frac1N\mu_k^2|\hat d_k|^2$. Conjugate modes $k$ and $N-k$
contribute equally, so $S=c_0+2\sum_{k=1}^{11}c_k+c_{12}$. The near-Nyquist half carries
about $97.6\%$ of $S$ for this orbit; the constant and low-frequency modes carry the
remaining $2.4\%$.

\begin{table}[h!]
\centering
\begin{tabular}{rrrr}
\toprule
$k$ & $16\sin^4(\pi k/24)$ & $|\hat d_k|^2/N$ & $c_k$\\
\midrule
0  & 0.0000  & 570.375 & 0.000\\
1  & 0.0046  & 17.919  & 0.083\\
2  & 0.0718  & 3.375   & 0.242\\
3  & 0.3431  & 19.125  & 6.563\\
4  & 1.0000  & 3.375   & 3.375\\
5  & 2.1974  & 2.331   & 5.122\\
6  & 4.0000  & 3.375   & 13.500\\
7  & 6.3385  & 2.331   & 14.774\\
8  & 9.0000  & 3.375   & 30.375\\
9  & 11.6569 & 19.125  & 222.937\\
10 & 13.9282 & 3.375   & 47.008\\
11 & 15.4595 & 17.919  & 277.022\\
12 & 16.0000 & 3.375   & 54.000\\
\bottomrule
\end{tabular}
\caption{Spectral decomposition of the $(1,1)$ orbit's curvature. Conjugate modes $k$ and
$24-k$ contribute equally, so $S=c_0+2\sum_{k=1}^{11}c_k+c_{12}=1296$. The contribution
rises with frequency; the dominant pairs are $k=11,13$ ($277.0$ each) and $k=9,15$
($222.9$ each), with the Nyquist mode $k=12$ contributing $54.0$. The near-Nyquist half
($\sin^2(\pi k/24)\ge\frac12$, i.e.\ $6\le k\le 18$) accounts for about $97.6\%$ of $S$.}
\label{tab:spectrum}
\end{table}

\begin{rem}[The maximizer is the highest-frequency class]
The three modulo-$9$ classes are ordered by how much of their curvature sits in the
near-Nyquist half: it is $97.6\%$ of $S$ for the $1296$ class, against $94.0\%$ and
$86.6\%$ for the $972$ and $648$ classes. So among the three full-period orbit classes
modulo~$9$, the Fibonacci class is the one whose curvature is most concentrated near the
Nyquist frequency. This is a diagnosis, not a derivation: because $S$ \emph{is} the
high-pass energy, the statement reformulates maximality spectrally but does not explain
why the $(1,1)$ seed produces the most near-Nyquist spectrum in the first place
(see Section~\ref{sec:open}).
\end{rem}

\subsection{The half-period asymmetry, corrected}\label{sec:halves}

The identity $F_{n+12}\equiv -F_n\pmod 9$ (from $24=2\cdot 12$~\cite{Wall1960}) suggests
the second differences should satisfy $\kappa_{n+12}=-\kappa_n$, which would force the two
half-period curvature sums to be equal. \emph{This is false}, and an earlier version of
this note relied on it in error. The relation $\kappa_{n+12}=-\kappa_n$ holds for
$n=1,\dots,9$ but breaks at $n=0,10,11$---exactly the indices adjacent to the displayed
zeros at positions $11$ and $23$, where $d(0)=9$ severs the antisymmetry. Consequently the
halves are \emph{unequal}:
\[
\sum_{n=0}^{11}\kappa_n^2=459,\qquad \sum_{n=12}^{23}\kappa_n^2=837,\qquad 459+837=1296,
\]
and the entire imbalance $837-459=378$ comes from those three zero-crossing index pairs.
The factorization $1296=2\cdot 648$ is therefore a coincidence, not a consequence of
half-period symmetry; the curvature does not split evenly across the period.

\section{How Far Maximality Extends}\label{sec:maximality}

It is natural to ask whether the Fibonacci seed maximizes digital curvature for every
modulus. It does not. Define, for any modulus $m$ with even Pisano period, $S_{\max}(m)$ as
the maximum of $S$ over full-period orbits, computed with the same display convention.

\begin{description}
\item[Small moduli.] Among moduli $m\le 14$ with even Pisano period, $(1,1)$ attains
$S_{\max}(m)$ in every case; the first failure is $m=15$.
\item[Erosion with scale.] Over the $118$ moduli $m\le 120$ with even Pisano period,
$(1,1)$ attains $S_{\max}(m)$ for $78$ of them and fails for $40$. The first three failures
are $m=15,33,40$ (and these are the only failures below $41$); the full failure list
continues
\[
45,49,52,53,56,59,60,64,65,67,68,71,74,80,81,83,85,86,88,90,\dots,120 .
\]
The proportion of failures broadly increases with $m$ (reaching $60$--$70\%$ among moduli
in the $80$s--$100$s), though not monotonically.
\item[No arithmetic predictor.] No simple arithmetic rule we tested separates the failures
from the successes. In particular ``$5\mid m$'' fails badly in both directions ($33$ is a
failure with $5\nmid 33$; $20,25,30,35$ are multiples of $5$ for which $(1,1)$ still wins).
\end{description}

The spectral lens supplies the robust statement that survives the failures. In every
exceptional case we examined---including $m=15,33,40$---the \emph{actual} maximizer is the
orbit carrying the most near-Nyquist (high-frequency) energy, and it carries strictly more
of it than the $(1,1)$ orbit does. For example at $m=40$ the maximizer carries $93.7\%$ of
its curvature in the near-Nyquist half versus $88.8\%$ for $(1,1)$. The invariant content of
``Fibonacci is special'' is thus spectral: among small moduli the classical cycle happens
to be the most high-frequency orbit, but for larger moduli other orbits are more jagged
still, and curvature tracks jaggedness rather than the Fibonacci seed.

\section{What Does Not Hold}\label{sec:negatives}

We record two negative results, both retiring claims from earlier versions, because closing
off false directions is part of the contribution.

\begin{description}
\item[The closed form $S_{\max}=m(\pi(m)/2)^2$ is a coincidence.] It holds at exactly
$m=9$ ($1296$) and $m=27$ ($34992$) and nowhere else among the moduli tested. It fails
already at $m=3$ (actual $36$, predicted $48$), and on the next powers of three at $m=81$
(actual $857304$, predicted $944784$) and $m=243$ (actual $21677544$, predicted
$25509168$); it never holds off the powers of three. The appealing reading
$1296=9F_{12}=9\cdot 144$ is numerology specific to $m=9$. There is no closed form for
$S_{\max}(m)$ supported by the data.
\item[The phenomenon does not generalize across recurrences.] Replacing the Fibonacci
recurrence by $a_{n+1}=p\,a_n+q\,a_{n-1}$ destroys ``simplest seed maximizes curvature''.
For $(p,q)=(1,2)$ neither $(0,1)$ nor $(1,1)$ ever maximizes (over the moduli tested); for
the Tribonacci recurrence the natural simple seeds almost never maximize. Only the
generator seed $(0,1)$ under $(p,q)=(2,1)$ shows a comparable, weaker effect. The clean
behaviour is special to Fibonacci.
\end{description}

\section{Computational Verification}\label{sec:verify}

The script \texttt{papers/v4/reproduce.py} in the project repository reproduces every
numerical claim in this note, asserting the exact ones: the trichotomy
$\{648{:}24,\,972{:}24,\,1296{:}24\}$ of Theorem~\ref{thm:trichotomy}, the spectral
identity of Proposition~\ref{prop:spectral}, the unit-class scalings of
Remark~\ref{rem:units}, the half-period split of Section~\ref{sec:halves}, the maximality
sweep of Section~\ref{sec:maximality}, and the failure of the closed form in
Section~\ref{sec:negatives}. It uses only the Python standard library. The original
\texttt{verify\_curvature.py}, which accompanies Versions~1--3, independently enumerates
all $81$ seeds modulo~$9$ and checks the trichotomy. Every numerical claim in this note has
thus been verified by at least two independent implementations.

\section{Revision History and Errata}\label{sec:errata}

This is Version~4. It supersedes Versions~1--3
(\href{https://doi.org/10.5281/zenodo.15811767}{doi:10.5281/zenodo.15811767}), which
remain available as the historical record. The substantive corrections are:
\begin{itemize}
\item \textbf{The trichotomy is proved by enumeration.} The congruence/``scaling''
argument offered as a proof in Versions~1--3 is vacuous (all three curvature values are
divisible by~$9$) and has been replaced by the finite verification, which is the actual
proof. The valid residue of that argument---the three values as the three classes of
$(\mathbb Z/9\mathbb Z)^\times/\{\pm1\}$---is kept as Remark~\ref{rem:units}.
\item \textbf{The half-period symmetry claim was wrong.} Version~2 asserted
$S=2\sum_{n=0}^{11}\kappa_n^2$ from an antisymmetry that the display convention breaks; the
true halves are $459$ and $837$ (Section~\ref{sec:halves}). Version~3 already withdrew the
explicit formula; this version explains precisely why it fails.
\item \textbf{The conjecture $S_{\max}=m(\pi(m)/2)^2$ is false}
(Section~\ref{sec:negatives}); it is a two-point coincidence at $m=9,27$.
\item \textbf{Section~6's ``explanation'' is replaced} by the spectral identity
(Section~\ref{sec:spectral}), which is the genuine reason the maximum is what it is.
\item \textbf{The novelty claim is softened.} We claim only that we are not aware of prior
work applying this spectral/Dirichlet-energy lens to the digital-root orbits of a Pisano
period; the underlying circulant diagonalization is classical. A search of OEIS
(A030132 and neighbours) and of MathSciNet for ``digital root'' and ``Pisano period''
turned up no prior treatment of these second-difference statistics, but we make no
exhaustive-priority claim.
\end{itemize}

\section{Open Questions}\label{sec:open}

\begin{itemize}
\item Characterize the moduli for which $(1,1)$ fails to maximize curvature
(Section~\ref{sec:maximality}); no arithmetic predictor was found, and roughly a third of
the failures have a maximizer that is \emph{not} a unit multiple of the Fibonacci orbit.
\item Make the spectral heuristic rigorous: prove (or refute) that the curvature
maximizer always maximizes the $\mu_k^2$-weighted near-Nyquist mass.
\item Is the $S_{\max}=m(\pi(m)/2)^2$ coincidence at $m=9,27$ part of any larger pattern,
or genuinely isolated?
\end{itemize}

\section*{Disclosure statement}
\noindent The author reports no competing interests.

\begin{thebibliography}{9}

\bibitem{Wall1960}
D.~D.~Wall,
``Fibonacci Series Modulo $m$,''
\textit{Amer.\ Math.\ Monthly} \textbf{67} (1960), 525--532.

\bibitem{OEIS}
OEIS~A030132, ``Digital root of Fibonacci($n$).''

\bibitem{Koshy2018}
T.~Koshy,
\textit{Fibonacci and Lucas Numbers with Applications},
Wiley (2018).

\bibitem{Davis1979}
P.~J.~Davis,
\textit{Circulant Matrices},
Wiley (1979).

\end{thebibliography}

\end{document}
